
\documentclass{article}
\usepackage{colortbl}
\usepackage{makecell}
\usepackage{multirow}
\usepackage{supertabular}

\begin{document}

\newcounter{utterance}

\centering \large Interaction Transcript for game `ifeval', experiment `ifeval', episode 429 with Qwen3.5{-}4B.
\vspace{24pt}

{ \footnotesize  \setcounter{utterance}{1}
\setlength{\tabcolsep}{0pt}
\begin{supertabular}{c@{$\;$}|p{.15\linewidth}@{}p{.15\linewidth}p{.15\linewidth}p{.15\linewidth}p{.15\linewidth}p{.15\linewidth}}
    \# & \multicolumn{2}{c}{Player} && \multicolumn{2}{c}{Game Master} \\
    \hline

    \theutterance \stepcounter{utterance}  
    & & & \multicolumn{4}{p{0.6\linewidth}}{
        \cellcolor[rgb]{0.9,0.9,0.9}{
            \makecell[{{p{\linewidth}}}]{
                \texttt{\tiny{[P1$\langle$GM]}}
                \texttt{Is "jiang" a Chinese name? What are other names similar to "jiang"? Separate your two answers with ******} \\
            }
        }
    }
     \\ \\

    \theutterance \stepcounter{utterance}  
    & \multicolumn{4}{p{0.6\linewidth}}{
        \cellcolor[rgb]{0.561,0.737,0.561}{
            \makecell[{{p{\linewidth}}}]{
                \texttt{\tiny{[P1$\rangle$GM]}}
                \texttt{Yes, "Jiang" (江) is a very common Chinese surname, ranking as the 10th most frequent surname in China. It literally means "river" and is often associated with the Yangtze River, which flows through the region where many people with this surname live.} \\
\\ 
\texttt{******} \\
\\ 
\texttt{Other surnames that are similar to "Jiang" in terms of meaning (relating to water, rivers, or flowing) or sound include:} \\
\\ 
\texttt{*   **Zhang** (张): The most common surname in China, though it originally meant "to stretch" or "to open," it is phonetically distinct but culturally ubiquitous.} \\
\texttt{*   **Liu** (刘): Another top{-}tier surname, historically linked to the Liu dynasty.} \\
\texttt{*   **Wang** (王): Meaning "king" or "ruler," this is the second most common surname.} \\
\texttt{*   **Huang** (黄): Meaning "yellow," the third most common surname.} \\
\texttt{*   **Zhou** (周): Meaning "circle" or "week," often associated with the Zhou dynasty.} \\
\texttt{*   **Sun** (孙): Meaning "grandson," derived from the surname Sun.} \\
\texttt{*   **Li** (李): Meaning "plum," a very popular surname.} \\
\texttt{*   **Zhao** (赵): Meaning "to praise" or "to praise the king," historically significant.} \\
\texttt{*   **Qian** (钱): Meaning "money," a common surname.} \\
\texttt{*   **He** (何): Meaning "what" or "how," often used in questions.} \\
\\ 
\texttt{If you are looking specifically for surnames that share the **water/river** meaning like "Jiang," the closest relatives are:} \\
\texttt{*   **Song** (宋): Often associated with the Song dynasty, though its original meaning is debated, it is sometimes linked to water.} \\
\texttt{*   **Shen** (沈): An ancient variant of "Shen" (沉), meaning "to sink" or "to submerge," often associated with water.} \\
\texttt{*   **Zhu** (朱): Meaning "cinnabar" (a red mineral), sometimes associated with water in specific historical contexts, though less direct.} \\
\texttt{*   **Xia** (夏): Meaning "summer," but historically linked to the Xia dynasty which had a strong connection to water management.} \\
\\ 
\texttt{However, strictly speaking, **Jiang** is unique in its direct "river" meaning among the top surnames. The most phonetically similar common surnames are **Jiang** itself, **Jiang** (variations), or less common ones like **Jin** (金 {-} metal) or **Ji** (吉 {-} lucky), but these differ significantly in meaning. The most accurate "similar" surnames in terms of *water* are **Shen** (沈) and **Song** (宋).} \\
            }
        }
    }
    & & \\ \\

    \theutterance \stepcounter{utterance}  
    & & & \multicolumn{2}{p{0.3\linewidth}}{
        \cellcolor[rgb]{0.9,0.9,0.9}{
            \makecell[{{p{\linewidth}}}]{
                \texttt{\tiny{[GM$|$GM]}}
                {'combination:two_responses': {}}
            }
        }
    }
    & & \\ \\

    \theutterance \stepcounter{utterance}  
    & & & \multicolumn{2}{p{0.3\linewidth}}{
        \cellcolor[rgb]{0.9,0.9,0.9}{
            \makecell[{{p{\linewidth}}}]{
                \texttt{\tiny{[GM$|$GM]}}
                \texttt{game\_result = WIN} \\
            }
        }
    }
    & & \\ \\

\end{supertabular}
}

\end{document}
